\chapter{Creatinine Blood Test and Kidney Function}

\section*{Overview}
A creatinine blood test measures creatinine, a muscle-waste product cleared by the kidneys. It is used with estimated glomerular filtration rate (eGFR) to assess kidney function.

\section*{Why This Test or Procedure Is Ordered}
It is commonly used to screen for kidney disease, monitor known kidney impairment, assess dehydration or systemic illness, and guide medication dosing.

\section*{How This Test or Procedure Is Performed}
A blood sample is collected from a vein. A urine creatinine measurement or urine albumin-to-creatinine ratio may be ordered alongside it.

\section*{Results and Limitations}
Results are interpreted with eGFR, age, body composition, hydration, and prior values. Muscle mass, diet, pregnancy, and some medicines can affect creatinine independently of kidney injury.

\section*{References}
\begin{enumerate}
  \item MedlinePlus. Creatinine Test. U.S. National Library of Medicine; 2025.
  \item Current Medical Diagnosis and Treatment. McGraw-Hill; 2025.
\end{enumerate}
\clearpage
